\section{Experimental Setup}
\label{sec:setup}

\paragraph{Three evaluation suites.}
Each suite answers a different experimental question. The \emph{eight-game main
timing/regression suite} in \cref{tab:suite} contains six public games via OpenSpiel
\cite{lanctot2019openspiel} and two native HUNL-style subgames
\cite{ganzfried2015endgame,brown2017safe} (\cref{app:hunl}). The suite spans 54 to
275{,}983 Infosets and supports the headline timing, operation-count, and
reference-parity results. The \emph{twelve-game update-rule suite} adds four more
OpenSpiel games for the four-rule, two-order CPU study of \cref{tab:convergence}. The
\emph{extended poker suite} contains the released Libratus river endgames \cite{brown2018superhuman,steinberger2019pokerrl} and larger
exported poker instances of the head-to-head in \cref{app:libratus}.

\input{figures/TABLE_T1_suite.tex}
% The ladder/scaling figure* is sourced here (not in sec 7.2 where it is
% discussed) so the double-column float can land on the same page as the
% suite table; a double float cannot enter a page that has already begun.


\paragraph{Backends.}
\emph{GPU-CFR} runs the compiled float32 solver on one NVIDIA A100 80GB in eager or
CUDA Graph mode. The two modes differ only in dispatch, and a vanilla CFR arm isolates
the engine from regret matching$^{+}$. \emph{Ours CPU} runs the same compiled tensor dataflow on eight threads
pinned to dedicated physical cores, with parallelism supplied by PyTorch's
intra-operator pool \cite{paszke2019pytorch}. \emph{LiteEFG} 0.1.5 \cite{liu2024liteefg} uses a single-threaded C++ backend. The
main matrix uses its simultaneous vanilla CFR mode, and \cref{app:extra} reports its
alternating CFR$^{+}$ preset. \emph{Kim (2026)}~\cite{kim2026parallelizing},
released as \texttt{noregret}, executes sequence-form
\cite{koller1997representations} CFR through CuPy sparse products on the same A100. The main matrix matches its update order
to GPU-CFR's simultaneous updates and reports its GPU backend. \emph{OpenSpiel} \cite{lanctot2019openspiel} contributes its Python CFR
implementation and exact best response. Every backend reports exploitability as NashConv$/2$. \Cref{tab:semantics} lists
the rule, update order, averaging, precision, and execution mode per arm.

\paragraph{Shared game representations.}
Public games are instantiated through OpenSpiel and converted to each backend's input
format. For the HUNL-style subgames and Libratus endgames, we export the compiled
intermediate representation to LiteEFG's generic extensive-form text format and load
the same checksummed tree through the Kim adapter. Structural counts and small-game behavior are checked across conversions
(\cref{app:protocol}), and \cref{app:libratus} reports the poker value-agreement
certificate.

\paragraph{Protocol.}
Timing rows run 1{,}000 CFR$^{+}$ iterations. OpenSpiel runs 200 CFR iterations and
is reported per iteration. GPU steady-state measurements use CUDA events \cite{guide2020cuda} after 50
warmup iterations. Graph capture and tree construction are measured separately, while
training-call rows include all work after solver construction. Each cell runs in an
independent low-load process, and tables report
medians over repeated runs. \Cref{app:protocol,app:correctness} gives the software
stack, repetition counts, load controls, and run-level traceability.
